\section{Introduction}
\label{sec:intro}

Audio-language models answer natural-language questions about soundscapes by
combining sound recognition with language-based reasoning
\cite{qwen25omni2025,funaudiochat2026}. Benchmarks such as the Massive Multi-Task
Audio Understanding and Reasoning benchmark span speech, environmental sounds,
and music \cite{mmau2024}. A central question is whether an answer is genuinely
supported by the acoustic evidence. Correct responses can often be recovered
from text alone \cite{authorship2026impact}, and improvements in audio question
answering can partly reflect better text-based reasoning \cite{omnir1}. For
environmental sounds this dependence is also local, since a question about a
brief bark concerns one event within a longer recording. Evaluating the answer
therefore requires linking it to the class of that event.

Intervention methods examine evidence use by perturbing the audio and observing
how the answer moves. Noise and region-guided masking probe audio faithfulness
\cite{lalmfaith2026}. Attention-guided masking with a duration-matched control
connects localized perturbations to prediction changes \cite{mapo2026}.
Counterfactual voice grounding tests speaker-attributed reasoning
\cite{hear2026whosaidwhat}, and paralinguistic audits vary acoustic cues while
holding transcripts fixed \cite{paraling2026}. These approaches provide useful
tests of acoustic grounding. For a local environmental-sound question, an answer
change depends on what the intervention changes at the queried location. A
replacement can change the event's class or rewrite the same interval with
another recording of that class. The matched same-class replacement measures
responses to local rewriting while preserving the queried class. Comparing it
with a different-class replacement asks how much additional answer movement
comes from changing that class under the same editing procedure. Removal and
masking tests instead probe what happens when local evidence is deleted or
degraded. The same
concern arises in text-domain counterfactual prompting, where a targeted edit
is a compound treatment that benefits from a meaning-preserving baseline
\cite{yang2026compared}.

This paper addresses how to attribute the movement of an answer after a local
waveform intervention. \emph{When an audio-language model changes its answer
after a local waveform intervention, how can we tell whether it responded to
the event's class or to the fact that the waveform was edited?} We introduce a
matched local waveform intervention protocol on synthetic soundscapes with
annotated event intervals, and make three contributions.
(i)~\textbf{A matched replacement protocol}. Class-changing and class-preserving
replacements share the target span, donor-fitting procedure, level estimate,
and crossfade. The class-preserving arm measures response to local waveform
rewriting. Nuisance, background-window, silence, and sham controls complete the
design.
(ii)~\textbf{A measure of answer changes and their direction}. We report
whether the answer changes and whether it identifies the replacement class.
An exact decomposition relates the observed gap to an accuracy-based
independence reference and covariance between condition outcomes.
(iii)~\textbf{A cross-architecture behavioral result}. Across Qwen2.5-Omni-7B
and Qwen2.5-Omni-3B \cite{qwen25omni2025} and Fun-Audio-Chat-8B
\cite{funaudiochat2026}, class-changing replacements produce 49.7 to 60.9
percentage points more answer changes than matched class-preserving rewrites.
Sham answer changes stay below 1.0\%, and soundscape-weighted selective
accuracy reaches 73.7 to 81.9\%.
